% Chapter generated from Part~III, Draft~10 by booktools/assemble.py
\begingroup
% macros of Part~III, Draft~10
\providecommand{\A}{}\renewcommand{\A}{\mathbb A}
\providecommand{\C}{}\renewcommand{\C}{\mathbb C}
\providecommand{\Z}{}\renewcommand{\Z}{\mathbb Z}
\providecommand{\GL}{}\renewcommand{\GL}{\operatorname{GL}}
\providecommand{\Nrd}{}\renewcommand{\Nrd}{\operatorname{Nrd}}
\providecommand{\Irr}{}\renewcommand{\Irr}{\operatorname{Irr}}
\providecommand{\vol}{}\renewcommand{\vol}{\operatorname{vol}}
\providecommand{\tH}{}\renewcommand{\tH}{\widetilde H}
\providecommand{\tZ}{}\renewcommand{\tZ}{\widetilde Z}
\providecommand{\ttT}{}\renewcommand{\ttT}{\widetilde T}
\providecommand{\reg}{}\renewcommand{\reg}{\mathrm{reg}}
\providecommand{\gen}{}\renewcommand{\gen}{\mathrm{gen}}
\providecommand{\JL}{}\renewcommand{\JL}{\mathrm{JL}}
\providecommand{\Sha}{}\renewcommand{\Sha}{\operatorname{Sha}}
\providecommand{\Gb}{}\renewcommand{\Gb}{\overline{G}}
\providecommand{\GDb}{}\renewcommand{\GDb}{\overline{G}_D}
\providecommand{\Zb}{}\renewcommand{\Zb}{\overline{Z}}
\providecommand{\Ptag}{}\renewcommand{\Ptag}{\textnormal{\textbf{(P)}}}
\providecommand{\Etag}{}\renewcommand{\Etag}{\textnormal{\textbf{(E)}}}
\providecommand{\Otag}{}\renewcommand{\Otag}{\textnormal{\textbf{(O)}}}
\providecommand{\Ctag}{}\renewcommand{\Ctag}{\textnormal{\textbf{(C)}}}
\chapter[The division-algebra side]{The division-algebra side and its local character transform}\label{ch:III}
\paragraph{Overview.} Let $D/F$ be a central division algebra of prime degree $r$ and
$\widetilde{D^\times}$ the $n$-fold cover ($n$ odd) specified in
Chapter~\ref{ch:I}, where $(n,r)=1$. We construct and normalize a unitary
transform $L_D$ of genuine and ordinary invariant class functions,
with local coefficient $|n|_F^{-(r-1)/2}$. When the residue characteristic
is prime to $n$, a unique splitting over the reduced-norm-one group,
together with an Adams operation, proves that $L_D$ bijects irreducible
characters and preserves dimensions. At places dividing $n$ we obtain
a nonsplitting and divisibility obstruction. We give the normalized
orbital-integral transfer and its trace-pairing identity. Globally we
prove an almost-everywhere local-to-global principle for $n$th powers
modulo scalars, and explain the relevant rational-orbit and measure
considerations from Kottwitz's work. The tame character relation with the cover of $\GL_r$ follows, within
its published correspondence range, by an explicit cocycle, star-map,
central-character and root-sum check. The wild-place character
correspondence and the adelic trace-formula comparison remain open.

\section{Introduction and results}
The classical local Jacquet--Langlands correspondence identifies
irreducible representations of $D^\times$ with essentially
square-integrable representations of $\GL_r(F)$ in the appropriate
central-character families \cite{DKV84}. It does not produce a
correspondence between ordinary representations of $D^\times$ and
genuine representations of its metaplectic cover: pullback from the
linear group is not genuine. This chapter studies the division-algebra
side directly and provides local results independent of a conjectural
global trace-formula comparison.

This section places the chapter in the series, records the status of its results, and introduces the running example. It depends only on the statements of Chapters~\ref{ch:I} and~\ref{ch:II}; references to them are to Chapter~\ref{ch:I}, and Chapter~\ref{ch:II}. The notation of the series and of this chapter is collected in \S\ref{III:sec:glossary}, the results are listed with their labels in the \hyperref[III:sec:summary]{Summary of results} that follows this section, and related work is discussed in \S\ref{III:sec:related}.

\paragraph{Position within the series.}
Chapter~\ref{ch:I} constructs the specified covers, computes their good elements, and
proves matching for regular-supported functions; Chapter~\ref{ch:II} proves matching of
arbitrary functions (its Theorem~E), identifies the germs at the center, and
proves an improved global transfer (its Theorem~C$'$). Here the central object
is the transform $L_D$ on the compact-mod-center group $D^\times$.
The principal issues and their status are:
\begin{center}
\small
\begin{tabular}{@{}p{0.23\linewidth}p{0.57\linewidth}p{0.10\linewidth}@{}}
\toprule
\textbf{Topic} & \textbf{Assertion} & \textbf{Status}\\
\midrule
Covering theory & Scalar centrality and good-locus description from Chapter~\ref{ch:I} & (E)\\
Local harmonic analysis & Unitary $L_D$ and explicit coefficient $b_D$ & (P)\\
Tame representations & Splitting over $D^1$ and irreducible-character bijection for $p\nmid n$ & (P)\\
Wild representations & Nonsplitting and necessary divisibility constraints & (P)\\
 & Integrality of $L_D$ at places dividing $n$ & (C)\\
Character relation (D2) of Chapter~\ref{ch:I} & Tame places, with the Flicker--Kazhdan--Zou theorem available (in particular when $(n,N_r)=1$); cocycle and star-map compatibility checked below (Corollary~\ref{III:cor:D2tame}) & (P)\\
Test functions & Normalized local orbital-integral transfer and trace adjunction & (P)\\
Global arithmetic & Local-to-global powers modulo scalars outside any finite set & (P)\\
Global comparison & Good-element matching, adelic measures, spectral separation & (C)\\
\bottomrule
\end{tabular}
\end{center}
The labels are those of the whole series (\S\ref{gl:labels}). As in Chapters~\ref{ch:I} and~\ref{ch:II}, \textbf{(P)} means proved here (possibly using
results marked \textbf{(E)}), \textbf{(E)} means established elsewhere
(including Chapters~\ref{ch:I} and~\ref{ch:II}) and used as a black box, \textbf{(O)} means that
a proof is outlined but not fully written out or checked, and \textbf{(C)}
means conjectural: new mathematical input is required.
These statuses concern the precise stated hypotheses; the local
claims rely on the specific covering cocycle of Chapter~\ref{ch:I}, not on
an arbitrary central extension.

The two distinct maps in the chapter should be kept separate:
the unitary operator on invariant $L^2$ class functions, and a
bijection of individual irreducible characters. The first exists at
all division places, whereas the second is established here only
when $p\nmid n$. The Adams operation provides the bridge in that case.

After the local results, we review the part of Kottwitz's arithmetic
framework needed for rational conjugacy and inner forms, then prove
a Hasse principle for the global power quotient. The final section
states what remains for an adelic metaplectic Jacquet--Langlands
comparison. Section~\ref{III:sec:glossary} collects the notation and glossary, and Section~\ref{III:sec:related} discusses related work; the literature of the whole series is reviewed in \S\ref{I:sec:lit}.

\subsection*{The running example}
In the series' running example, $K=\mathbb{Q}(\zeta_3)$, $n=3$, $r=2$, $D=(-1,4-\zeta_3)_K$, ramified exactly at $S=\{\mathfrak p_3,\mathfrak p_7\}$, with $\mathfrak p_3=(1-\zeta_3)$ and $\mathfrak p_7=(3+\zeta_3)$ (Example~\ref{ex:running}, and \cite{JRE26}), the two places of $S$ illustrate the two regimes of this chapter. At $\mathfrak p_7$, which does not divide $n$, $b_D=c_D=1$, the discriminant ratio is identically $1$ (Lemma~\ref{III:lem:tame-R}), and $L_D$ is a dimension-preserving bijection of irreducible characters (Theorem~\ref{III:thm:tame-transfer}). At $\mathfrak p_3$, where $K_v=\mathbb{Q}_3(\zeta_3)$ and $|3|_v=1/9$, one has $b_D=3$ and $c_D=27$. There the cover does not split, and genuine irreducibles have dimension divisible by $3$ (Proposition~\ref{III:prop:wild-obstruction}). Integrality of $L_D$ would force $L_D\mathbf 1$ to be irreducible of dimension $27$ (Remark~\ref{III:rem:wild-test}). Whether it is remains open; see Chapter~\ref{ch:IV}.

\section{Local setting and good elements}\label{III:sec:setting}
We fix the local field, central-character conventions, and the
covering-theoretic theorem imported from Chapter~\ref{ch:I}. The distinction
between a good element of the cover and an ordinary regular element
is essential to the definition of the transform. The section depends on \S\S\ref{I:sec:covers}--\ref{I:sec:tori} (the cover, scalar centrality, good elements). It introduces $F$, $D$, $H$, $Z$, $\tH$, $\tZ$, the characters $\tilde\theta$ and $\theta$, the star map, the discriminant $\Delta$ and good regular elements.
Let $F$ be a finite extension of $\mathbb Q_p$ containing $\mu_n$, where $n\ge1$ is odd. Let $r\ge2$ be prime and $(n,r)=1$. Let $D/F$ be a central division algebra of degree $r$, $H=D^\times$, and $Z=F^\times$ its scalar center. We use the normalized local absolute value $|\cdot|_F$ ($|\varpi_F|_F=q_F^{-1}$).

Let
\begin{equation}\label{III:eq:extension}
1\longrightarrow\mu_n\longrightarrow\tH\overset{p}{\longrightarrow}H\longrightarrow1
\end{equation}
be a locally compact topological central extension. Throughout we use \emph{the specified cocycle of Chapter~\ref{ch:I}}, rather than an arbitrary central extension. In its parameter convention, $n\mid r-1+2rc_v$; under that condition the inverse image $\tZ=p^{-1}(Z)$ is central (\cref{I:lem:center}). The cocycle condition is essential; a general $n$-fold extension need not have central scalar preimage. References to Chapter~\ref{ch:I} below refer to its stated cocycle conventions.

In the notation of Chapters~\ref{ch:I} and~\ref{ch:II}, $\tH=\GDb(F)$, $\tZ=\Zb$, and $\Gb_v$ denotes the Kazhdan--Patterson cover of $\GL_r(F)$ with parameter $c_v$; the normalized orbital integral $J=\Delta O$ of \S\ref{III:sec:orbital} is the function $F=\Delta\Phi$ of \cref{I:def:oi}, for the measures used here. Fix a unitary character $\widetilde\theta:\tZ\to\C^\times$ whose restriction to $\mu_n$ is faithful. Write $\theta(a)=\widetilde\theta(a^*)$ for $a\in Z$, where the star map is defined next. For invariant functions we impose the central equivariance conventions $f(ah)=\theta(a)f(h)$ and $\widetilde f(zh)=\widetilde\theta(z)\widetilde f(h)$.

\begin{definition}[The star map]\label{III:def:star}\index{star map@Star map}
For $h\in H$ choose any lift $\widetilde h\in\tH$ and set
\begin{equation}\label{III:eq:star}
h^*=\widetilde h^{\,n}\in\tH.
\end{equation}
The definition is independent of the lift. It is conjugation equivariant, and for $a\in Z$ and $h\in H$ one has $(ah)^*=a^*h^*$.
\end{definition}
\begin{proof}
Another lift is $\zeta\widetilde h$ with $\zeta\in\mu_n$, whose $n$th power is $\widetilde h^n$. Conjugation equivariance follows by lifting the conjugating element. Elements of $\tZ$ are central by assumption, hence $(\widetilde a\widetilde h)^n=\widetilde a^n\widetilde h^n$. In particular $a\mapsto a^*$ is a homomorphism from $Z$ to $\tZ$. (For a maximal torus it is also a homomorphism, because $n$ is odd; see below.)
\end{proof}

Every noncentral $h\in D^\times$ generates a maximal subfield $E=F(h)$, of degree $r$. Set $T=E^\times$; it is its own centralizer in $H$. A regular element of $H$ means a noncentral element. The corresponding Weyl discriminant is
\begin{equation}\label{III:eq:disc}
\Delta(h)=\left|\det\bigl(1-\operatorname{Ad}(h)\mid \mathfrak h/\mathfrak t\bigr)\right|_F^{1/2}
=\prod_{i<j}\left|\frac{(\lambda_i-\lambda_j)^2}{\lambda_i\lambda_j}\right|_F^{1/2},
\end{equation}
where $\lambda_1,\dots,\lambda_r$ are the conjugate eigenvalues in a splitting field, and the product is interpreted with the normalized absolute value.

\begin{definition}[Good regular element]\label{III:def:good}\index{good regular element@Good regular element}
A regular $\widetilde h\in\tH$ is good if it commutes with $p^{-1}(Z_H(p(\widetilde h)))$. The definition is equivalent to the usual statement that the commutator character on the preimage of the regular centralizer is trivial. Every $h^*$ is good, since its commutators are $n$th powers in $\mu_n$.
\end{definition}

\begin{remark}\label{III:rem:good-vanish}
If $\widetilde h$ is regular and not good, every genuine conjugation-invariant function $\widetilde f$ vanishes at $\widetilde h$: some lift $\widetilde t$ centralizing $p(\widetilde h)$ has $\widetilde t\widetilde h\widetilde t^{-1}=\zeta\widetilde h$, with $\zeta\neq1$; invariance and genuine central equivariance then force $\widetilde f(\widetilde h)=0$.
\end{remark}

\begin{proposition}[\Etag\ Good lifts on the division side (Chapter~\ref{ch:I})]\label{III:prop:good}
For each maximal subfield $E/F$ of $D$, the good elements of
$p^{-1}(E^\times)$ are precisely $z\,t^*$ with $z\in\tZ$ and
$t\in E^\times$. Equivalently, modulo $\tZ$ the good subset
corresponds to $(E^\times/F^\times)^n$. These are the good-element
and commutator results established in \cref{I:lem:center} and
\cref{I:prop:good}(a); the local index and preservation of regularity
are treated there in \cref{I:prop:good}(b),(d).
\end{proposition}

\section{Power maps on elliptic tori}\label{III:sec:power}
The local index of the power map and the ratio of Weyl discriminants
control the normalization of $L_D$. Both calculations use the
coprimality of $n$ and $r$, including at wild places. The section depends on \S\ref{III:sec:setting} and introduces the discriminant ratio $\Delta(t)/\Delta(t^n)$.
\begin{lemma}[No new singular elements]\label{III:lem:regular}
If $h\in H$ is noncentral, $h^n$ is noncentral. In particular the image of a regular element under the star map is regular.
\end{lemma}
\begin{proof}
Suppose $h^n\in F^\times$. The field $L=F(h)$ is an intermediate field of a maximal subfield of $D$, hence $[L:F]$ divides $r$. Because $F$ contains $\mu_n$ and $h^n\in F^\times$, Kummer theory shows $[L:F]$ divides $n$. The coprimality $(n,r)=1$ therefore gives $[L:F]=1$, contrary to the assumption that $h$ is noncentral.
\end{proof}

\begin{remark}[Why $(n,r)=1$ cannot be dropped]\label{III:rem:coprimality-needed}
The proof isolates exactly where coprimality is used: $[L:F]$ must divide both $r$ (it sits inside a degree-$r$ extension) and $n$ (Kummer theory, since $h^n$ is a scalar), and $(n,r)=1$ is what forces the common divisor down to $1$. Drop the hypothesis and the lemma is false, by an explicit example at a wild place of the running example's own field.

Take $F=K_{\p_3}=\Q_3(\zeta_3)$, the local field of \S\ref{sec:running}, but now with $n=r=3$ instead of $n=3,r=2$ (so $(n,r)=3\ne1$). Fix a central division algebra $D/F$ of degree $3$ (one exists over every local field, for every degree). A uniformizer $\varpi$ of $F$ has valuation $v(\varpi)=1$, not divisible by $3$, so $\varpi$ is never a cube in $F^\times$; hence $L=F(\varpi^{1/3})$ is a genuine, totally ramified, cyclic cubic Kummer extension of $F$ (cyclic because $\mu_3\subset F$), and it embeds as a maximal subfield of $D$. Its generator $h=\varpi^{1/3}$ is noncentral, since $[F(h):F]=3\ne1$, but
\[
h^n=h^3=\varpi\in F^\times,
\]
a scalar. So the star map sends the regular element $h$ to a central, hence singular, element: the exact failure the lemma rules out when $(n,r)=1$.

This is the same arithmetic as the elementary fact that ``multiply by $n$'' on $\Z/r\Z$ is a bijection exactly when $\gcd(n,r)=1$: here the power map on $E^\times/F^\times$ for a degree-$r$ field $E$ plays the role of that multiplication, and losing injectivity on the subgroup generated by $h$ is losing it on all of $E^\times/F^\times$ by Lemma~\ref{III:lem:index} below. The failure is not cosmetic: Definition~\ref{III:def:L} of $L_D$ is only written down on the regular set, exactly because the formula needs $t^*=t^n$ to again be regular; at $n=r$, the computation above produces a regular $t$ whose star $t^*$ is central, so $L_D$ as defined would have to assign a value to $zt^*$ coming from every regular class with that same central image, which is inconsistent unless those classes already agree. The construction of \S\ref{III:sec:local} simply does not get off the ground without $(n,r)=1$.
\end{remark}

\begin{lemma}[Degree of the power map]\label{III:lem:index}
For any degree-$r$ extension $E/F$ embedded in $D$, the map
\[
P_n:E^\times/F^\times\longrightarrow E^\times/F^\times,\qquad [t]\longmapsto[t^n]
\]
is injective and has open image of index
\begin{equation}\label{III:eq:index}
I_F(n,r)=\bigl[E^\times/F^\times:(E^\times/F^\times)^n\bigr]
=|n|_F^{-(r-1)}.
\end{equation}
The index is independent of $E$.
\end{lemma}
\begin{proof}
If $t^n\in F^\times$, Kummer theory over a field containing $\mu_n$ gives $[F(t):F]\mid n$. Since $F(t)\subset E$ and $[E:F]=r$, this degree also divides $r$; hence $t\in F^\times$, proving injectivity.

The quotient $A=E^\times/F^\times$ is a compact abelian $p$-adic Lie group. Its $p$-adic analytic dimension is $d=([E:\mathbb Q_p]-[F:\mathbb Q_p])=(r-1)[F:\mathbb Q_p]$. The structure theorem for compact abelian $p$-adic Lie groups writes $A\simeq \mathbb Z_p^d\times C$ for a finite abelian group $C$. The injectivity already proved implies $A[n]=\{1\}$, so multiplication by $n$ is an automorphism of $C$. On $\mathbb Z_p^d$ it has index $p^{v_p(n)d}$. Since $|n|_F=p^{-[F:\mathbb Q_p]v_p(n)}$, this is exactly $|n|_F^{-(r-1)}$. (The argument also handles $p\mid n$.)
\end{proof}

\begin{lemma}[Discriminant ratio]\label{III:lem:ratio}
For noncentral $h$, the ratio
\[
R(h)=\frac{\Delta(h)}{\Delta(h^n)}
\]
is locally constant and nowhere zero. It extends uniquely to a locally constant function on all of $H$, taking the value
\begin{equation}\label{III:eq:ratio-at-center}
R(a)=|n|_F^{-r(r-1)/2}\qquad(a\in F^\times).
\end{equation}
\end{lemma}
\begin{proof}
By Lemma~\ref{III:lem:regular}, both discriminants are nonzero at every noncentral element. They are continuous functions with values in the discrete multiplicative group $q_F^{\frac12\mathbb Z}$ and hence locally constant on this open set.

For $h$ sufficiently close to a scalar $a\in F^\times$, all its eigenvalues are close to $a$ in a finite splitting extension. Factoring
\[
\lambda_i^n-\lambda_j^n=(\lambda_i-\lambda_j)
\sum_{k=0}^{n-1}\lambda_i^{n-1-k}\lambda_j^k,
\]
the final factor approaches $na^{n-1}$ and thus has constant absolute value $|na^{n-1}|$ for $h$ in a sufficiently small neighborhood. The factors involving $\lambda_i\lambda_j$ in the normalized discriminant cancel the $|a|^{n-1}$ powers. Each of the $r(r-1)/2$ unordered pairs consequently contributes $|n|_F^{-1}$ to $R(h)$. This proves the formula and locally constant extension; uniqueness follows from density of the regular set.
\end{proof}

\section{The normalized local class-function transform}\label{III:sec:local}
We first transfer \emph{invariant functions}; the problem of sending
individual irreducible characters to irreducible characters is
addressed separately. The transform is normalized using quotient
probability measures. The section depends on \S\S\ref{III:sec:setting}--\ref{III:sec:power}. It introduces the probability measures, the Hilbert spaces $\mathcal H_\theta(H)$ and $\mathcal H_{\widetilde\theta}(\tH)$, the transform $L_D$ and the constant $b_D$.
We equip $H/Z$ with Haar probability measure. Because $\tZ=p^{-1}(Z)$, the quotient $\tH/\tZ$ maps isomorphically to $H/Z$; we transport the probability measure through this isomorphism. Each elliptic torus quotient $T/Z$ likewise has its probability Haar measure. These conventions are the relevant ones for \emph{local unitarity}; matching geometric distributions globally may require additional, separately stated orbital-integral conventions.

Write $\mathcal H_\theta(H)$ for the Hilbert space of $L^2$ conjugation-invariant functions on $H$ with $\theta$-equivariance under $Z$. Similarly, write $\mathcal H_{\widetilde\theta}(\tH)$ for the genuine invariant functions with $\widetilde\theta$-equivariance under $\tZ$. The inner products are taken over $H/Z$ and $\tH/\tZ$, respectively.

\begin{definition}[The division-side transform]\label{III:def:L}\index{division side transform@Division-side transform, $L_D$}
Using Proposition~\ref{III:prop:good}, define on regular elements
\begin{equation}\label{III:eq:transfer}
\boxed{\quad
(L_D\varphi)(z\,t^*)=
|n|_F^{-(r-1)/2}\,\widetilde\theta(z)
\frac{\Delta(t)}{\Delta(t^n)}\,\varphi(t),
\quad z\in\tZ,
\ t\in T^{\reg},\quad}
\end{equation}
and let $L_D\varphi=0$ at nongood regular elements. The value at singular elements is understood by continuous extension whenever $\varphi$ is smooth.
\end{definition}

\begin{lemma}[Well-definedness and smoothness]\label{III:lem:well-defined}
The expression \eqref{III:eq:transfer} is independent of the choice of $z,t$ and of the maximal subfield $E=F(t)$. It sends locally constant invariant functions to locally constant invariant functions.
\end{lemma}
\begin{proof}
For a good regular element, the regular centralizer is uniquely $E^\times$. If $zt^*=z'(t')^*$ then $t'^n$ and $t^n$ differ by a scalar. Since their centralizers are $E$ and $t',t\in E^\times$, the injectivity in Lemma~\ref{III:lem:index} gives $t'=at$ with $a\in F^\times$. As $(at)^*=a^*t^*$ and $\Delta(at)=\Delta(t)$, the formula changes by the same factor $\theta(a)=\widetilde\theta(a^*)$ on either side. Conjugation equivariance follows from Definition~\ref{III:def:star}.

On a regular semisimple conjugacy chart, the image of the $n$th-power map on $T/Z$ is an open-and-closed subgroup by Lemma~\ref{III:lem:index}. Thus the prescribed function is locally constant on both the good and nongood regular loci, using Lemma~\ref{III:lem:ratio}. A sufficiently small neighborhood of any central element consists of elements represented by $z t^*$: the local $n$th-root map exists by the nonarchimedean inverse function theorem, since $n$ is nonzero in $F$. The same discriminant ratio has the constant value of \eqref{III:eq:ratio-at-center} there. Hence the function extends locally constantly across the center. Since $r$ is prime there are no other nonregular elements.
\end{proof}

\begin{theorem}[Unitarity and explicit constant]\label{III:thm:unitary}
For the specified cover of Chapter~\ref{ch:I}, $L_D$ extends uniquely to a unitary isomorphism
\begin{equation}\label{III:eq:unitary}
L_D:\mathcal H_\theta(H)\xrightarrow{\ \sim\ }\mathcal H_{\widetilde\theta}(\tH).
\end{equation}
The positive coefficient $b_D=|n|_F^{-(r-1)/2}$ in \eqref{III:eq:transfer} is uniquely determined by unitarity for the quotient probability measures used here.
\end{theorem}
\begin{proof}
The Weyl integration formula, applied to $|\varphi|^2$, reduces the norm on $H/Z$ to a sum over elliptic maximal tori $T/Z$ of weighted integrals
\[
\frac{1}{|W(H,T)|}\int_{T/Z}\Delta(t)^2|\varphi(t)|^2\,dt.
\]
There is no need for a separate Weyl integration theorem on the covering group: since $\tH/\tZ\simeq H/Z$, the function $|\widetilde f|^2$ is an honest conjugation-invariant function on $H/Z$, and the ordinary Weyl formula applies directly. By Remark~\ref{III:rem:good-vanish} and Proposition~\ref{III:prop:good}, a genuine invariant function vanishes off the open-and-closed subset $(T/Z)^n$ of each torus quotient.

The power map $P_n:T/Z\to(T/Z)^n$ is a topological group isomorphism. Since $T/Z$ has probability Haar measure and $(T/Z)^n$ has index $I=|n|_F^{-(r-1)}$ in $T/Z$, for every integrable function $Q$ on this subgroup,
\begin{equation}\label{III:eq:changevar}
\int_{(T/Z)^n}Q(x)\,dx=\frac1I\int_{T/Z}Q(t^n)\,dt.
\end{equation}
Consequently, with a provisional coefficient $b>0$, the covered Weyl integral equals
\[
\frac{b^2}{I\,|W(H,T)|}\int_{T/Z}\Delta(t)^2|\varphi(t)|^2\,dt.
\]
The index is independent of $T$, so the equality of norms for all class functions is equivalent to $b^2=I$. This gives $b=b_D$.

To see surjectivity, start with a smooth genuine invariant function on the cover. It vanishes off the good regular set. By the uniqueness in Lemma~\ref{III:lem:well-defined}, the formula defining $L_D$ can be inverted on every good regular element. The inverse has the same local smoothness properties, by the inverse of $P_n$ on its open image. This proves surjectivity on a dense subspace and therefore the unitary isomorphism on completions.
\end{proof}

\begin{remark}[Local versus global normalization; comparison with Zou]\label{III:rem:normalizations}
The coefficient $b_D=|n|_F^{-(r-1)/2}$ is determined by the
power-map Jacobian and the quotient probability measures in
Theorem~\ref{III:thm:unitary}. A common compatible rescaling of Haar
measures does not alter the ratio, but the proposed global trace
formula requires independent, explicit adelic normalizations.

Zou's corrected local Flicker--Kazhdan character relation
\cite[(3.2) and Remark~3.3]{Z23} has elliptic coefficient
\[
  \frac{n}{d_F(d_r)|n|_F^{r/2}},\qquad
  d_F(m)=m|m|_F^{-1/2},\qquad
  d_r=\gcd(n,2rc+r-1).
\]
Under Chapter~\ref{ch:I}'s scalar-centrality congruence one has $d_r=n$, and
this coefficient simplifies algebraically to $b_D$.
To identify the \emph{character transforms}, however, one must also
identify the two cocycles, their star maps, the central-character
conventions, and the discriminant and orbital-measure normalizations.
The equality of numerical coefficients alone is not such an
identification; see Corollary~\ref{III:cor:D2tame}.
\end{remark}

\section{Irreducible characters and the integrality obstruction}\label{III:sec:characters}
The unitary transform has a matrix relative to irreducible characters,
but unitarity alone does not force that matrix to be integral.
We record explicit arithmetic tests that any true character
correspondence must pass. The section depends on \S\ref{III:sec:local} and introduces the character matrix of $L_D$ and the constant $c_D$.
Since $H/Z$ is compact, every smooth irreducible representation with a fixed central character is finite dimensional. Its character is a locally constant invariant function on $H$. The same holds on the cover.

\begin{corollary}[A unitary character matrix]\label{III:cor:matrix}
Let $\rho$ range over $\Irr_\theta(H)$ and $\tau$ over $\Irr_{\widetilde\theta}(\tH)$. For each $\rho$ there is a \emph{finite} expansion
\begin{equation}\label{III:eq:matrix}
L_D\Theta_\rho=\sum_{\tau}a_{\rho\tau}\Theta_\tau.
\end{equation}
The matrix $A=(a_{\rho\tau})$ is unitary (in the sense of the associated operator on $\ell^2$), and its rows and columns are finite.
\end{corollary}
\begin{proof}
Schur orthogonality for compact groups modulo the fixed center gives orthonormal bases of irreducible characters for the two invariant $L^2$ spaces. Theorem~\ref{III:thm:unitary} therefore makes $A$ unitary. Local constancy and compactness modulo center show that each smooth invariant function is invariant under an open normal subgroup of the compact quotient, so it factors through a finite quotient (or through a corresponding finite-dimensional projective quotient with fixed central character). Its irreducible-character expansion is therefore finite. Applying this observation to the inverse transform gives finite columns as well.
\end{proof}

\begin{corollary}[What integral coefficients would imply]\label{III:cor:integrality}
Suppose the coefficients $a_{\rho\tau}$ in \eqref{III:eq:matrix} are integers for every $\rho,\tau$. Then $L_D$ maps each irreducible character to a unique irreducible character with positive sign, and it induces a bijection
\[
\Irr_\theta(H)\;\xrightarrow{\sim}\;\Irr_{\widetilde\theta}(\tH).
\]
More precisely, if $L_D\Theta_\rho=\Theta_\tau$, then
\begin{equation}\label{III:eq:dimension}
\boxed{\qquad\dim\tau=|n|_F^{-(r^2-1)/2}\,\dim\rho.\qquad}
\end{equation}
\end{corollary}
\begin{proof}
An integral vector of squared $\ell^2$-norm one has exactly one nonzero entry, equal to $+1$ or $-1$. The sign cannot be negative because $\Theta_\rho(1)=\dim\rho>0$ and the continuous value at the identity prescribed by \eqref{III:eq:transfer} is
\[
(L_D\Theta_\rho)(1)
=b_D\,|n|_F^{-r(r-1)/2}\dim\rho>0.
\]
The same argument applied to $L_D^{-1}$ yields bijectivity. Substituting $b_D=|n|_F^{-(r-1)/2}$ gives \eqref{III:eq:dimension}.
\end{proof}

\begin{remark}[A concrete wild-place check]
Take $r=2$, $n=3$, $F=\mathbb Q_3(\zeta_3)$. Then $|3|_F=3^{-2}=1/9$, so $I=9$, $b_D=3$, and formula \eqref{III:eq:dimension} predicts $\dim\tau=27\dim\rho$ \emph{if} integral irreducible transfer holds. This is a sharp, falsifiable numerical consequence of the proposed correspondence; it is not evidence that the correspondence is already proved.
\end{remark}


\begin{lemma}[Integral local scaling constants]\label{III:lem:integral-constants}
If $p\mid n$ and $\mu_n\subset F$, with $n$ odd as assumed, then
\[
b_D=|n|_F^{-(r-1)/2},\qquad
c_D=|n|_F^{-(r^2-1)/2}
\]
are positive integers.
\end{lemma}
\begin{proof}
Write $p^k\Vert n$ and $d=[F:\mathbb Q_p]$. Since
$\mu_{p^k}\subset F$ and $p$ is odd, $p-1$ divides $d$;
in particular $d$ is even. For the normalized absolute value,
$|n|_F^{-1}=p^{kd}$ is a perfect square. Both exponents above
therefore give integral nonnegative powers of $p$.
\end{proof}

\begin{remark}[A falsifiable wild-place divisibility test]\label{III:rem:wild-test}
Suppose the matrix of Corollary~\ref{III:cor:matrix} is integral at a
wild place, so the transform is a bijection on irreducible characters.
Every genuine $\tau\in\Irr_{\widetilde\theta}(\tH)$ must then satisfy
$c_D\mid\dim\tau$, since $\dim\tau=c_D\dim\rho$ for its unique
preimage $\rho$. In the example $(r,n,F)=(2,3,\mathbb Q_3(\zeta_3))$,
$c_D=27$, while the covered-torus argument of
Proposition~\ref{III:prop:wild-obstruction} yields only the weaker
factor $b_D=3$ in the dimension constraints. A genuine irreducible
of dimension not divisible by $27$ would disprove \emph{integrality
of this specific transform} $L_D$, but would not disprove the existence
of some other wild-place correspondence. The existence of such an
irreducible is an open computational test, not an asserted example.
\end{remark}

\begin{remark}[Why unitarity is not integrality]\label{III:rem:warning}
A unitary operator need not send an orthonormal basis to itself. Even if its matrix is locally finite, the entries need not be integral. A trace-formula or representation-theoretic positivity argument is needed to show $a_{\rho\tau}\in\mathbb Z$. The remaining genuinely unresolved local step is the integrality statement at \emph{wild places}, where $p\mid n$; see Section~\ref{III:sec:wild}. It should not be concealed under the phrase ``unitary character transfer.''
\end{remark}

\section{Adams operations and finite-group characters}\label{III:sec:adams}
We recall the $K$-theoretic origin of the Adams operations and
explain why they preserve irreducibility in certain tame
finite-group situations. The section depends on nothing earlier in the chapter and introduces the Adams operations $\psi^n$ and virtual characters.
\subsection*{Background from $K$-theory and finite-group characters}
The Adams operations (after J.~F.~Adams) originated in algebraic topology, where complex topological
$K$-theory studies the Grothendieck group of complex vector bundles. For a finite
group $Q$, a $Q$-equivariant complex vector bundle over a point is simply a
finite-dimensional complex representation of $Q$. Consequently
\[
K_Q^0(\mathrm{pt})\cong R(Q),
\]
the representation (or character) ring \cite{Seg68}. The operations
$\lambda^j[V]=[\bigwedge^jV]$ give this ring its $\lambda$-ring structure.
The Adams operations $\psi^k$ are the corresponding Newton-polynomial
operations; for instance
\[
\psi^1(V)=V,\qquad
\psi^2(V)=V\otimes V-2\bigwedge\nolimits^2V
\quad\text{in }R(Q).
\]
This description connects Adams's topological construction \cite{Ada62}
with Atiyah's systematic account of power operations \cite{Ati66}.

To see the character-theoretic formula, restrict $V$ to the cyclic subgroup
generated by $g\in Q$. On this subgroup $V$ decomposes as a sum of eigenlines
with eigenvalues $\alpha_i$; the defining property $\psi^k(L)=L^{\otimes k}$
on line representations gives
\begin{equation}\label{III:eq:adams-def-background}
\boxed{\quad\chi_{\psi^k(V)}(g)=\sum_i\alpha_i^k
      =\chi_V(g^k).\quad}
\end{equation}
The operation is integral on the Grothendieck ring, but need not preserve the
positive cone of actual characters. For example, if $V$ is the standard
two-dimensional irreducible representation of $S_3$, then
$\psi^2(V)=\mathbf1-\operatorname{sgn}+V$, a virtual character.
If $\gcd(k,|Q|)=1$, however, raising the root-of-unity eigenvalues to the
$k$th power is a cyclotomic Galois automorphism; it permutes irreducible
characters. Section~\ref{III:sec:tame} exploits a more delicate version of this
fact with a \emph{fixed central character}, so that a central $n$-primary
subgroup can be retained even when the full finite quotient has order
divisible by $n$. This is why the Adams argument establishes integrality at
tame places but not at wild ones.

For a genuine finite-dimensional representation $\tau$ of $\tH$, let
\[
\Psi_n(\Theta_\tau)(h)=\Theta_\tau(h^*)=\operatorname{tr}\tau(\widetilde h^n).
\]
The $n$th Adams operation in the complex representation ring of $\tH$ is a virtual character with integer coefficients. Since $\mu_n$ acts on $\tau$ by scalars, its $n$th Adams operation is trivial on $\mu_n$ and descends to a virtual character of $H$. Thus $\Psi_n$ is integral and has central character $\theta$.

\begin{proposition}[Adams pullback and its weight]\label{III:prop:adams}
For $\tau\in\Irr_{\widetilde\theta}(\tH)$, the class function $\Psi_n(\Theta_\tau)$ is a virtual character in the integral representation ring $R_\theta(H)$. On regular $t\in H$,
\begin{equation}\label{III:eq:adamsweight}
\Psi_n(\Theta_\tau)(t)=
|n|_F^{-(r-1)/2}\,R(t)\,\bigl(L_D^{-1}\Theta_\tau\bigr)(t).
\end{equation}
\end{proposition}
\begin{proof}
After a suitable unramified twist, a given finite-dimensional smooth representation has finite image, and hence factors through a finite quotient $Q$ of $\tH$. The usual $n$th Adams operation is a virtual character of $Q$. For $\zeta\in\mu_n$ it satisfies $\psi^n\Theta(\zeta g)=\Theta((\zeta g)^n)=\Theta(g^n)$, so linear independence of irreducible characters shows that all its constituents have trivial $\mu_n$-action. Thus it descends to $H$, with the required central character. Undoing the unramified twist preserves this assertion (it changes under the $n$th power in the expected way). The displayed formula is the inverse of \eqref{III:eq:transfer} at $z=1$.
\end{proof}

In general the nonconstant weight $R(t)^{-1}$ prevents us from deducing integrality of $L_D$ from integrality of the Adams operation alone. In Section~\ref{III:sec:tame}, however, we prove $R\equiv1$ when $p\nmid n$, and then obtain a complete tame-case character correspondence. At wild places, the discriminant weight is generally nontrivial.


\section{The tame case: splitting and irreducible character integrality}\label{III:sec:tame}
In this section assume that the residue characteristic $p$ does not divide $n$.
We use the special cover of Chapter~\ref{ch:I}, including the centrality of $\tZ$ and
its good-element calculation. Write $H^1=\ker(\Nrd:H\to F^\times)$,
and put $U_D^1=1+\mathfrak P_D$ for the principal units in the maximal
order of $D$. The term \emph{tame} here means $p\nmid n$; it does not
impose a condition on the ramification of a maximal subfield of $D$. The section depends on \S\S\ref{III:sec:setting}--\ref{III:sec:adams}, on Chapter~\ref{ch:I}, and, for Corollary~\ref{III:cor:D2tame}, on \cite{DKV84} and \cite{Z23}. It introduces $H^1$ and $U_D^1$.

\begin{lemma}[Prime-to-$n$ finite quotients]\label{III:lem:prime-to-n}
All finite continuous quotients of $H^1$ and $H/Z$ have order relatively
prime to $n$. In particular both groups have continuous cohomology
$H^i(H^1,\mu_n)=0$ for $i=1,2$, where $H^1$ acts trivially on $\mu_n$.
\end{lemma}
\begin{proof}
Write $q=|\mathcal O_F/\mathfrak p_F|$. The standard residue-field
quotient of the norm-one group embeds in
\[
\ker\bigl(N_{\mathbb F_{q^r}/\mathbb F_q}:
\mathbb F_{q^r}^\times\longrightarrow\mathbb F_q^\times\bigr),
\qquad M:=\frac{q^r-1}{q-1};
\]
the kernel of reduction is pro-$p$. In fact the displayed finite group
is the full residue quotient, but divisibility of its order by $M$ is
all that is needed. Since $\mu_n\subset F$ and $p\nmid n$,
$q\equiv1\pmod n$, and therefore $M\equiv r\pmod n$.
Consequently $\gcd(M,n)=1$ by $\gcd(r,n)=1$, and no finite quotient
of $H^1$ can have order divisible by a prime factor of $n$.

For $H/Z$, the image of $U_D^1$ is pro-$p$, and the residual quotient
has an extension with one factor of order dividing $M$ and another
of order $r$ (the valuation modulo scalar valuations). Thus every
finite quotient of $H/Z$ has order prime to $n$ as well.
For a profinite group all of whose finite quotients have order
prime to $n$, continuous cohomology with trivial finite
$\mu_n$-coefficients vanishes in positive degrees: every cocycle
factors through a finite quotient, where averaging is available.
\end{proof}

\begin{proposition}[Unique splitting over the norm-one group]\label{III:prop:tame-split}
The restriction of the cover to $H^1$ admits a unique continuous
homomorphic splitting $s:H^1\to\tH$. Its image is normal in $\tH$.
The quotient
\[
\mathcal A=\tH/s(H^1)
\]
is an abelian locally compact group, fitting into a central extension
$1\to\mu_n\to\mathcal A\to F^\times\to1$.
For each prescribed faithful character $\epsilon$ of $\mu_n$ there exists
\emph{some} continuous unitary character $\nu:\tH\to\mathbb C^\times$
with $\nu|_{\mu_n}=\epsilon$ and $\nu|_{s(H^1)}=1$.
\end{proposition}
\begin{proof}
Vanishing of $H^2(H^1,\mu_n)$ in Lemma~\ref{III:lem:prime-to-n}
provides a continuous splitting, while vanishing of $H^1$
gives uniqueness. Conjugating a splitting by a lift of an element
of $H$ produces another splitting of the normal group $H^1$;
uniqueness therefore proves normality.

The reduced norm identifies $H/H^1$ with $F^\times$.
The commutator of two elements of $\mathcal A$ gives a continuous
alternating bicharacter
$c:F^\times\times F^\times\to\mu_n$.
The preimage of scalar elements is central in $\tH$, and the
reduced norm of a scalar $a$ is $a^r$. Hence $c(a^r,b)=1$.
By bilinearity $c(a,b)^r=1$, while $c(a,b)^n=1$ since its image
lies in $\mu_n$. Coprimality of $r,n$ implies $c=1$, so
$\mathcal A$ is abelian. Every continuous unitary character of a
closed subgroup of a locally compact abelian group extends to
such a character of the whole group. Apply this to the finite
closed subgroup $\mu_n\subset\mathcal A$ and compose with
$\tH\to\mathcal A$.
\end{proof}

\begin{remark}[Twisting is not the character transfer]\label{III:rem:tame-twist}
A choice of $\nu$ produces an elementary dimension-preserving
bijection between irreducible ordinary representations and genuine
irreducible representations, via
$\rho\mapsto\nu\otimes(\rho\circ p)$.
This statement concerns the \emph{full} genuine dual, with central
characters adjusted by $\nu$. It is not, by itself, the assertion
that the canonical transform $L_D$ realizes this particular bijection.
Moreover a prescribed $\widetilde\theta$ need not equal
$\nu|_{\tZ}$; keeping track of this difference is necessary.
\end{remark}

\begin{lemma}[Tame discriminant identity]\label{III:lem:tame-R}
For every $h\in H\setminus Z$ one has
\[
\boxed{\;\Delta(h^n)=\Delta(h),\qquad R(h)=1.\;}
\]
The locally constant extension of $R$ to the scalar center is also
identically one.
\end{lemma}
\begin{proof}
The compact profinite group $K=H/Z$ has all finite quotients of
order prime to $n$ by Lemma~\ref{III:lem:prime-to-n}. On each finite
quotient $Q$ the power map $x\mapsto x^n$ is bijective (its inverse
is $x\mapsto x^m$ for any $m$ with $mn\equiv1\pmod{\exp Q}$).
Consequently $P_n:K\to K$ is a homeomorphism which preserves
probability Haar measure: it permutes the cosets of every open
normal subgroup. It is an analytic local diffeomorphism, since
its derivative is invertible at regular points by
Lemma~\ref{III:lem:regular}, and at scalars it is multiplication by
$n\ne0$ on the Lie algebra.

At a regular element $h$, write $h\exp(X)$ for a local
left-trivialized chart. The derivative of the power map on
$\mathfrak h/\mathfrak z$ is
\[
(dP_n)_h=\sum_{k=0}^{n-1}\operatorname{Ad}(h)^{-k}.
\]
On the $r-1$ dimensional subspace $\mathfrak t/\mathfrak z$,
$\operatorname{Ad}(h)$ acts trivially, so these directions contribute
$|n|_F^{r-1}$ to the absolute determinant. Over a splitting field,
the complementary root directions have eigenvalues
$u_{ij}=\lambda_i/\lambda_j$ for $i\ne j$ under
$\operatorname{Ad}(h)$. The corresponding eigenvalues of the
derivative are
\[
\sum_{k=0}^{n-1}u_{ij}^{-k}
=\frac{1-u_{ij}^{-n}}{1-u_{ij}^{-1}},
\]
which are nonzero by Lemma~\ref{III:lem:regular}. Taking the product
over $i\ne j$ (equivalently, the $F$-determinant) yields
\[
|\det(dP_n)_h|_F
=|n|_F^{r-1}\prod_{i\ne j}
  \left|\frac{1-u_{ij}^{-n}}{1-u_{ij}^{-1}}\right|_F
=|n|_F^{r-1}\frac{\Delta(h^n)^2}{\Delta(h)^2}.
\]
Here the middle product is interpreted with the appropriate
extension of the normalized absolute value to a splitting field.
The $p$-adic analytic change-of-variables theorem
\cite[Proposition~1.2.1]{Loe08} applies to the local analytic
diffeomorphism $P_n$. The probability Haar measure is associated
locally to a nonvanishing left-invariant top-degree differential
form, so preservation of Haar measure implies
$|\det(dP_n)_h|_F=1$ almost everywhere. The absolute determinant
is locally constant and nonzero on the regular set, hence is $1$
at every regular $h$. Since $|n|_F=1$, the discriminant identity
follows. The value at scalars is $1$ by Lemma~\ref{III:lem:ratio}.
\end{proof}

\begin{lemma}[Adams operations with fixed central character]\label{III:lem:tame-adams}
Take $G=D^\times$ and $Z_G=F^\times$, and assume $p\nmid n$.
For every smooth unitary character $\delta$ of $Z_G$, the
$n$th Adams operation gives a bijection
\[
\psi^n:\Irr_\delta(G)\xrightarrow{\sim}
\Irr_{\delta^n}(G),\qquad
\Theta_{\psi^n\sigma}(g)=\Theta_\sigma(g^n),
\]
with central characters interpreted on the same scalar center.
\end{lemma}
\begin{proof}
First consider a finite group $Q$ with central subgroup $C$
and $|Q/C|$ prime to $n$. Let $\pi$ be the set of primes dividing $n$, and let $Q_\pi$
be the $\pi$-primary subgroup of the finite abelian group $C$.
Since $|Q/C|$ is prime to $n$, we have
$|Q:Q_\pi|=|Q/C|\,|C:Q_\pi|$ prime to $n$.
Thus $Q_\pi$ is a \emph{normal Hall $\pi$-subgroup} of $Q$
(rather than a single-prime Sylow subgroup);
it is central because $Q_\pi\subset C$. Schur--Zassenhaus
provides a complement $Q_{\pi'}$, and centrality gives
$Q=Q_\pi\times Q_{\pi'}$, with $|Q_{\pi'}|$ prime to $n$.
Fixing the central character on $C$ fixes the one-dimensional
representation of $Q_\pi$.
On $Q_{\pi'}$ the operation $\psi^n$ is the cyclotomic Galois
automorphism $\zeta\mapsto\zeta^n$ of character values; since
$n$ is prime to $|Q_{\pi'}|$, it permutes irreducible characters.
On $Q_\pi$ it raises the fixed central character to its $n$th
power. Therefore $\psi^n$ bijects irreducibles of central
character $\delta$ with those of central character $\delta^n$.
This is a statement \emph{at a fixed $\delta$}: the operation
need not be injective on a union of different central-character
fibres, since distinct characters whose quotient has order dividing
$n$ can have the same $n$th power.

For $G=D^\times$, given $\delta$ choose a unitary unramified
character $\eta=\xi\circ\Nrd$ of $G$ so that
$\delta\eta|_Z$ is trivial on a chosen scalar uniformizer
$\varpi_F$ (choose an $r$th root of the inverse scalar value).
Then every smooth irreducible with central character $\delta\eta|_Z$
becomes a smooth representation of the compact profinite group
$G/\langle\varpi_F\rangle$; it factors through a finite quotient,
and its image central subgroup has quotient of order prime to
$n$. For the reverse direction, choose the open normal subgroup
through which a given representation factors small enough that
$\delta\eta|_Z$ is trivial on its intersection with the center;
this is possible because the center modulo $\langle\varpi_F\rangle$
is a closed subgroup of the compact profinite quotient.
Then both the specified central character and its $n$th power
factor through the same finite central subgroup. The preceding
finite-group argument applies to these representations.
Undo the twist, observing that
$\psi^n(\eta\otimes\sigma)=\eta^n\otimes\psi^n\sigma$.
The same finite-quotient argument applied in reverse proves
surjectivity. Note that this argument keeps the central
$\pi$-part of each finite quotient: an unramified twist cannot
generally remove $n$-torsion coming from the units of $Z$.
\end{proof}

\begin{theorem}[Tame irreducible character transfer]\label{III:thm:tame-transfer}
Assume $p\nmid n$ and the covering data of Chapter~\ref{ch:I}.
For each unitary genuine central character $\widetilde\theta$
on $\tZ$, put $\theta(a)=\widetilde\theta(a^*)$.
Then $L_D$ maps irreducible characters bijectively onto
genuine irreducible characters,
\[
\boxed{\quad
L_D:\{\Theta_\rho:\rho\in\Irr_\theta(H)\}
\xrightarrow{\sim}
\{\Theta_\tau:\tau\in\Irr_{\widetilde\theta}(\tH)\}.
\quad}
\]
Every representation on the right corresponding to $\rho$
has $\dim\tau=\dim\rho$. In particular the character matrix
of Corollary~\ref{III:cor:matrix} is a permutation matrix in the
tame case.
\end{theorem}
\begin{proof}
Choose the genuine one-dimensional character $\nu$ given by
Proposition~\ref{III:prop:tame-split}, with
$\nu|_{\mu_n}=\widetilde\theta|_{\mu_n}$. Put
\[
\chi_\nu(h)=\nu(\widetilde h)^n\quad(h\in H),\qquad
\delta(a)=\frac{\widetilde\theta(\widetilde a)}
{\nu(\widetilde a)}\quad(a\in F^\times).
\]
These are well-defined ordinary characters; in particular
$\chi_\nu$ is defined on all of $H$, while $\delta$ is a
character of $Z$. By the definition of the star map,
\[
\theta(a)=\chi_\nu(a)\,\delta(a)^n.
\]
For $\rho\in\Irr_\theta(H)$, the representation
$\chi_\nu^{-1}\otimes\rho$ has central character $\delta^n$.
By Lemma~\ref{III:lem:tame-adams} there is a unique
$\sigma\in\Irr_\delta(H)$ such that
\[
\Theta_\rho(t)=\chi_\nu(t)\,\Theta_\sigma(t^n)
\qquad(t\in H).
\]
Define the genuine irreducible representation
$\tau=\nu\otimes(\sigma\circ p)$.
It has central character $\widetilde\theta$, and for regular $t$
\[
\Theta_\tau(t^*)
=\nu(t^*)\,\Theta_\sigma(t^n)
=\chi_\nu(t)\,\Theta_\sigma(t^n)
=\Theta_\rho(t).
\]
By Lemma~\ref{III:lem:index} the power map is an automorphism of
each compact torus $E^\times/F^\times$ when $p\nmid n$.
Thus every regular element of $\tH$ is of the form $z t^*$
for some $z\in\tZ$ and regular $t$. The same equality holds
on all regular elements by central equivariance.
Since $b_D=1$ and $R=1$ by Lemma~\ref{III:lem:tame-R}, this is
exactly $L_D\Theta_\rho=\Theta_\tau$.
The map is bijective by the corresponding bijection for
$\psi^n$ and by tensoring with $\nu$. Finally tensoring with
one-dimensional characters and the finite-group Galois
operation preserve dimension.
\end{proof}

\begin{remark}[Independence of the auxiliary genuine character]
Although the construction in the proof chooses a genuine
one-dimensional character $\nu$, the resulting correspondence is
independent of that choice. Indeed, $L_D$ was defined solely from
$\widetilde\theta$, the star map, and the normalized discriminant.
For each $\rho$, the character $L_D\Theta_\rho$ determines a unique
irreducible genuine representation of $\tH$; this follows from
orthogonality (or linear independence) of irreducible characters.
Consequently any permissible choice of $\nu$ produces the same
isomorphism class $\tau$. This is a \emph{canonical correspondence
relative to the fixed covering datum and central characters}, not
an assertion that $\nu$ itself is canonical.
\end{remark}



\paragraph{Compatibility audit preceding Corollary~\ref{III:cor:D2tame}.}
We distinguish the intrinsically defined $n$th power of a lift from
any choice of cocycle section.  The following facts are the exact
compatibilities needed to compose the two local correspondences;
they are consequences of Chapter~\ref{ch:I} and the formulas in
\cite[\S\S2.2, 3.1]{Z23}, not additional assumptions.

Write $F=K_v$ and $c_v=r'c$ as in Chapter~\ref{ch:I}, with $rr'\equiv1\pmod n$.
\cref{I:prop:restriction}(b), identifies the restriction of the ambient
cover, after the prescribed relabeling of $\mu_n$, with the
Kazhdan--Patterson extension whose parameter is $c_v$.  In the
standard parameter-zero section the cocycle on the diagonal torus is
$\sigma_F^{(0)}(a,b)=\prod_{i<j}(a_i,b_j)_n$; adding the determinant
factor gives
\begin{equation}\label{III:eq:zou-cocycle-audit}
 \sigma_F^{(c_v)}(a,b)=
 \prod_{i<j}(a_i,b_j)_n
 \left(\prod_i a_i,\prod_j b_j\right)_n^{c_v}.
\end{equation}
This is precisely Zou's block formula \cite[(2.1)]{Z23} in the
one-dimensional-block case.  In particular the two models determine
the same extension class and the same centrality parameter
\[
 d_r=\gcd(n,r-1+2rc_v)=n.
\]
Here the claimed identification is of \emph{extensions}; Chapter~\ref{ch:I} does
not claim a preferred equality of two cocycles on all pairs of
nondiagonal matrices.

On scalar matrices $aI_r,bI_r$, equation
\eqref{III:eq:zou-cocycle-audit} becomes
\begin{equation}\label{III:eq:scalar-audit}
 \sigma_F^{(c_v)}(aI_r,bI_r)=(a,b)_n^{e_r},
 \qquad e_r=\binom r2+c_vr^2,
 \qquad 2e_r=r(r-1+2rc_v).
\end{equation}
Since $n$ is odd and divides $r-1+2rc_v$, it divides $e_r$.
Consequently the \emph{standard} Kazhdan--Patterson section is
multiplicative on scalar matrices, and for $a\in F^\times$ its
$n$th power satisfies $(s(aI_r))^n=s(a^n I_r)$.
In particular if $\zeta\in\mu_n(F)$, then $(\zeta I_r)^*=1$.
Thus our scalar character
$\theta(a)=\widetilde\theta(a^*)$ satisfies
$\theta|_{\mu_n}=1$ and is exactly Zou's required relation
$\widetilde\omega(s(a^nI_r))=\omega(a)$, after transporting
$\widetilde\theta$ through an isomorphism of covers.

Flicker--Kazhdan define the star map for odd $n$ by
$x^*=\widetilde x^{\,n}$ \cite[\S4]{FK86}; this agrees with
\eqref{III:eq:star} and is natural under \emph{every} isomorphism of
central extensions fixing $\mu_n$.  Two identifications of the
ambient $\GL_r$ cover with the Kazhdan--Patterson model differ
by $\chi\circ\det$ for a character $\chi:F^\times\to\mu_n$
(\cref{I:prop:restriction}(c)).  Their values agree on star images,
since $\chi(\det(x^n))=1$.  At arbitrary scalar translates one
transports the central character through the same identification.
Hence this choice does not alter the eventual identity for genuine
characters of the original ambient cover.

Finally, \cref{I:thm:ident}(b), identifies the covered elliptic
tori by conjugation inside the common ambient cover and is the
identity on their common central subgroup.  It therefore satisfies
\begin{equation}\label{III:eq:jstar-audit}
 \jmath(t_D^*)=(j(t))_{\GL}^*,\qquad
 \jmath(z t_D^*)=z\,(j(t))_{\GL}^*
 \quad(z\in\tZ).
\end{equation}
The Weyl discriminants on $D^\times$ and $\GL_r(F)$ are the same
expression in the conjugate eigenvalues of $t$; see
\eqref{III:eq:disc} and \cref{I:def:oi}.
These are character \emph{values}, not global orbital integrals:
no comparison of arbitrary quotient Haar measures is being inferred.

\begin{corollary}[Tame character relation in the available Flicker--Kazhdan range]\label{III:cor:D2tame}
Let $F=K_v$ be a finite place at which $D$ is division,
$p\nmid n$, and the standing odd-$n$ and $(n,r)=1$ assumptions
hold. Fix the common faithful genuine central character
$\widetilde\theta$ of Chapter~\ref{ch:I} and let
$\theta(a)=\widetilde\theta(a^*)$.
Assume that the square-integrable Flicker--Kazhdan correspondence
for the parameter-$c_v$ Kazhdan--Patterson cover is available
in the corrected form of \cite[Proposition~3.6 and
Definition~3.7]{Z23}. A sufficient published range for this
external input is $(n,N_r)=1$, where $N_r$ is the least common
multiple of the composite positive integers at most $r$, with
$N_r=1$ when there are none \cite[\S26]{FK86}.

Then the division-side tame transfer $L_D$ of
Theorem~\ref{III:thm:tame-transfer}, the classical local
Jacquet--Langlands correspondence, and the corrected metaplectic
correspondence compose to a bijection between the relevant
genuine irreducible representations $\widetilde\pi_D$ of
$\widetilde{D^\times}$ and genuine essentially square-integrable
representations $\widetilde\pi$ of $\Gb_v$ with the specified
central character.  On matching regular good elements,
\begin{equation}\label{III:eq:tame-D2}
 \boxed{\quad
 \Theta_{\widetilde\pi_D}(\widetilde\gamma)
 =(-1)^{r-1}\Theta_{\widetilde\pi}
       (\jmath(\widetilde\gamma)).\quad}
\end{equation}
Thus relation (D2) of Chapter~\ref{ch:I} is proved in this stated tame
range, conditional \emph{only} on the cited external
Flicker--Kazhdan theorem and the established local results of
Chapters~\ref{ch:I}--\ref{ch:III}; it is not asserted for all $(n,r)$.
\end{corollary}
\begin{proof}
\emph{1. Matching characters on the division side.}
By Theorem~\ref{III:thm:tame-transfer} there is a unique
$\rho\in\Irr_\theta(D^\times)$ for which
$L_D\Theta_\rho=\Theta_{\widetilde\pi_D}$.
At $p\nmid n$ we have $b_D=1$ and
$\Delta(t^n)=\Delta(t)$ (Lemma~\ref{III:lem:tame-R}); thus
\begin{equation}\label{III:eq:div-D2-step}
 \Theta_{\widetilde\pi_D}(z t_D^*)
   =\widetilde\theta(z)\Theta_\rho(t)
 \qquad(t\in E^\times\setminus F^\times,
            z\in\tZ).
\end{equation}
Every regular good element has this form (\cref{I:prop:good}(a)), and the power map is injective on
$E^\times/F^\times$ by Lemma~\ref{III:lem:index}.

\emph{2. Classical Jacquet--Langlands.}
Let $\pi=\JL(\rho)$, with the same ordinary scalar character
$\theta$.  The usual elliptic-character relation is
\begin{equation}\label{III:eq:linear-D2-step}
 \Theta_\rho(t)=(-1)^{r-1}\Theta_\pi(j(t)),
\end{equation}
where $j:E\hookrightarrow M_r(F)$ and the discriminant
normalization is the common one above \cite{DKV84}.
The character $\theta$ is trivial on $\mu_n$ by
\eqref{III:eq:scalar-audit}, so $\pi$ lies in the permitted
linear central-character family in Zou's correspondence.

\emph{3. Explicit evaluation of Zou's root sum.}
Let $x=j(t)$ be regular elliptic.  If $y^n=x^n$ in
$\GL_r(F)$, then $y$ commutes with $x^n$.  As $t^n$
is regular (Lemma~\ref{III:lem:regular}), its centralizer is
$j(E^\times)$; therefore $y\in j(E^\times)$ and
$(yx^{-1})^n=1$.  Since $\mu_n(E)=\mu_n(F)$, the
solutions are exactly the $n$ scalar multiples
$y=\zeta x$ with $\zeta\in\mu_n$.
For every one, $\Delta(y)=\Delta(x)$ and
$\Theta_\pi(y)=\Theta_\pi(x)$ because
$\theta(\zeta)=1$.  Equation~\eqref{III:eq:scalar-audit}
gives $(\zeta I_r)^*=1$, so the phase
$\epsilon(x^*/y^*)$ in Zou's equation (3.2) is also $1$.
Zou's corrected root-sum formula therefore reduces to
\begin{equation}\label{III:eq:zou-D2-step}
 \Theta_{\widetilde\pi}(x^*)
  =\frac{n}{d_F(d_r)|n|_F^{r/2}}
   \frac{\Delta(x)}{\Delta(x^n)}\Theta_\pi(x)
  =b_D\frac{\Delta(x)}{\Delta(x^n)}\Theta_\pi(x),
\end{equation}
where $d_F(m)=m|m|_F^{-1/2}$ and
$d_r=\gcd(n,r-1+2rc_v)=n$.  Indeed,
\[
 \frac{n}{d_F(n)|n|_F^{r/2}}
 =\frac{n}{n|n|_F^{-1/2}|n|_F^{r/2}}
 =|n|_F^{-(r-1)/2}=b_D.
\]
The equality of numerical factors is valid even at wild places;
it is \emph{not} a wild-place character-integrality theorem.
Here $p\nmid n$ implies $b_D=1$, and
Lemma~\ref{III:lem:tame-R} gives $\Delta(x)=\Delta(x^n)$.
Consequently $\Theta_{\widetilde\pi}(x^*)
=\Theta_\pi(x)$.

\emph{4. Transfer back to the two actual covers.}
By \eqref{III:eq:jstar-audit}, the star values in the last
identity coincide with those obtained from the covered
division-algebra torus.  Multiplying by
$\widetilde\theta(z)$ and combining
\eqref{III:eq:div-D2-step}, \eqref{III:eq:linear-D2-step}
and \eqref{III:eq:zou-D2-step} yields
\eqref{III:eq:tame-D2} for every good regular element.
The star and central-character identifications in the
preceding compatibility audit show that the conclusion is
independent of the auxiliary isomorphism to the standard
Kazhdan--Patterson cocycle model.

\emph{5. Bijection rather than just a character identity.}
Theorem~\ref{III:thm:tame-transfer} is a bijection on genuine
irreducibles with fixed central character; classical local
Jacquet--Langlands is a bijection onto the corresponding
essentially square-integrable linear representations; and
Zou's corrected correspondence is a bijection in the
specified external range and central-character family.
Their composition has an inverse obtained by reversing the
three maps.  This concludes the claim without inferring
irreducibility or integrality from $L^2$ unitarity alone.
\end{proof}
\begin{remark}[Scope and status of Corollary~\ref{III:cor:D2tame}]
The earlier version of this corollary carried status
\textbf{(O)} because the cocycle, star map and central
characters had not been compared explicitly with Zou.
Equations~\eqref{III:eq:zou-cocycle-audit}--\eqref{III:eq:jstar-audit}
now perform that comparison; the result is \textbf{(P)}
\emph{under the written external correspondence hypothesis}.
In particular $(n,N_r)=1$ suffices by the stated
Flicker--Kazhdan theorem, as used by Zou.  For additional
parameter pairs one must verify an applicable extension of
that external input; the present proof does not supply it.
At places $p\mid n$, integrality of $L_D$ remains open.
Global matching of Haar measures, test-function normalizations,
and unramified spherical units is also not proved by this
local pointwise character identity.
\end{remark}

\section{Places dividing \texorpdfstring{$n$}{n}: non-splitting and constraints}\label{III:sec:wild}
At residue characteristics dividing $n$, the tame splitting
argument fails. The commutator pairing from Chapter~\ref{ch:I} gives a precise
obstruction, but does not supply an irreducible-character
correspondence. The numerical test of Remark~\ref{III:rem:wild-test}
is an additional constraint on any proposed correspondence. The section depends on \S\S\ref{III:sec:setting} and~\ref{III:sec:characters} and on the commutator pairing of \S\ref{I:sec:tori}; it introduces no new notation.

\begin{proposition}[Wild-place obstruction to splitting]\label{III:prop:wild-obstruction}
Suppose $p\mid n$ and retain the commutator formula and
nondegeneracy assertion of \cref{I:prop:good}(a),(c).
Then $\tH$ does not split over $H^1$, and it has no genuine
one-dimensional representations. More precisely, every genuine irreducible
representation of $\tH$ has dimension divisible by
\[
 b_D=|n|_F^{-(r-1)/2}>1.
\]
The existence or integrality of the wild-place transform $L_D$
on irreducible characters is not established by this proposition.
\end{proposition}
\begin{proof}
Put $E^1=\ker(N_{E/F}:E^\times\to F^\times)$ and
$B=E^\times/(F^\times E^{\times n})$. The group $B$ has
order $|n|_F^{-(r-1)}>1$ by Lemma~\ref{III:lem:index}; it is
$n$-torsion. The map $E^1\to B$ is surjective: for $t\in E^\times$,
\[
 s_t=t^r/N_{E/F}(t)\in E^1,
\qquad [s_t]=[t]^r\in B.
\]
Since $\gcd(r,n)=1$, raising to the $r$th power is an
automorphism of $B$, so these elements generate $B$.
The commutator pairing on the covered torus, described in
\cref{I:prop:good}(a), has radical
$F^\times E^{\times n}$; it is therefore nondegenerate on
$B$. Restricting to $E^1$ preserves a nontrivial pairing
because $E^1\to B$ is surjective. Consequently the
inverse image of $E^1\subset H^1$ in $\tH$ is nonabelian,
so the cover cannot admit a splitting over $H^1$.

\cref{I:prop:good}(c), identifies the dimensions of
genuine irreducible representations of the covered torus
with $\sqrt{|B|}=b_D$ (for the fixed genuine central
character). Restricting a finite-dimensional genuine
representation of $\tH$ to the covered torus and decomposing
it into these constituents proves the divisibility assertion,
which excludes genuine one-dimensional representations.
\end{proof}

\begin{remark}[Global use and the remaining difficulty]\label{III:rem:global-tame}
At a finite place $v$ with $p_v\nmid n$, the local
$D_v^\times$-side integrality problem is now resolved under
the stated covering assumptions. At places dividing $n$
it remains a distinct and genuinely nonabelian problem.
Global division algebras can be chosen to avoid those places
subject to the Brauer-invariant sum and local-degree
conditions. This removes the \emph{division-side local
integrality obstacle} for such a choice, but does not prove
global trace-formula comparison, spectral multiplicity
statements, or a metaplectic Jacquet--Langlands theorem
at the split places.
\end{remark}

\section{Local orbital-integral transfer on \texorpdfstring{$D^\times$}{D x}}\label{III:sec:orbital}
This section depends on \S\S\ref{III:sec:setting} and~\ref{III:sec:local}. It introduces the spaces $\mathcal C_{\eta^{-1}}(H)$, the orbital integrals $O$ and $J=\Delta O$, the transforms $L_D^{(\widetilde\eta)}$, and the transfers $\mathcal O_D$ and $\mathcal B$. In this section no assumption $p\nmid n$ is imposed. The construction therefore
applies at both tame and wild places; it uses the good-element description of
Proposition~\ref{III:prop:good}. We keep the quotient probability Haar measures from
Section~\ref{III:sec:local}. The result is a transfer of test functions and geometric
distributions; it is not an assertion that the transform of every irreducible
character is irreducible.

\subsection{Test functions and orbital integrals}
For a unitary central character $\eta$ of $Z$, let
\[
\mathcal C_{\eta^{-1}}(H)
=\{f:H\to\mathbb C:\ f\text{ is locally constant and }
                   f(ah)=\eta(a)^{-1}f(h)\text{ for }a\in Z\}.
\]
We use analogous notation for $\mathcal C_{\widetilde\eta^{-1}}(\tH)$.
Since $H/Z$ and $\tH/\tZ$ are compact, all these functions are compactly
supported modulo the indicated center. These are the test-function central
characters \emph{inverse} to the central characters of representations:
$\Theta_\rho(ah)=\theta(a)\Theta_\rho(h)$, whereas
$f(ah)=\theta(a)^{-1}f(h)$, so that the product is a function on $H/Z$.

Write $dh$ for Haar probability measure on $H/Z$, and define the
conjugation average
\begin{equation}\label{III:eq:avg-H}
(\mathcal A_Hf)(h)
=\int_{H/Z}f(x^{-1}hx)\,dx.
\end{equation}
For the cover define, using any lift $\widetilde x$ of $x$,
\begin{equation}\label{III:eq:avg-cover}
(\mathcal A_{\tH}\widetilde f)(\widetilde h)
=\int_{H/Z}\widetilde f(\widetilde x^{-1}\widetilde h\widetilde x)\,dx.
\end{equation}
Centrality of $\mu_n$ and $\tZ$ makes the latter integrand independent of
the chosen lifts and of the representative of $x\in H/Z$. Both averages
are smooth, conjugation invariant, and central equivariant; they satisfy
$\mathcal A_H^2=\mathcal A_H$ and
$\mathcal A_{\tH}^2=\mathcal A_{\tH}$.

For $t\in T^{\reg}$, $T=F(t)^\times$, equip $H/T$ with the quotient
\emph{probability} measure induced by those on $H/Z$ and $T/Z$. Define
\begin{equation}\label{III:eq:ordinary-O}
O_H(t,f)=\int_{H/T}f(x^{-1}tx)\,d\dot x
           =(\mathcal A_Hf)(t).
\end{equation}
For a good regular $\widetilde t\in p^{-1}(T)$ the corresponding genuine
orbital integral is
\begin{equation}\label{III:eq:covered-O}
O_{\tH}(\widetilde t,\widetilde f)
=\int_{H/T}\widetilde f
  (\widetilde x^{-1}\widetilde t\widetilde x)\,d\dot x
=(\mathcal A_{\tH}\widetilde f)(\widetilde t).
\end{equation}
The first integral is independent of coset representatives precisely because
$\widetilde t$ is good. If $\widetilde t$ is not good, the full conjugation
average in \eqref{III:eq:avg-cover} vanishes: replacing $\widetilde x$ by
$\widetilde u\widetilde x$, for a lift of an element of $T$ with nontrivial
commutator against $\widetilde t$, multiplies the averaged value by a
nontrivial genuine scalar. We use this vanishing rather than attempting to
define an $H/T$ integral for a nongood element.

For later use put
\begin{equation}\label{III:eq:normalized-O}
J_H(t,f)=\Delta(t)O_H(t,f),\qquad
J_{\tH}(\widetilde t,\widetilde f)
=\Delta(p(\widetilde t))O_{\tH}(\widetilde t,\widetilde f)
\end{equation}
when $\widetilde t$ is good. The factor $\Delta$ is the \emph{square root}
of the ordinary Weyl discriminant as in \eqref{III:eq:disc}.

\subsection{An explicit transfer theorem}
Let $L_D^{(\widetilde\eta)}$ denote the map in
Definition~\ref{III:def:L} for the genuine central character
$\widetilde\eta$, with its associated ordinary central character
$\eta(a)=\widetilde\eta(a^*)$. In particular,
$L_D^{(\widetilde\theta^{-1})}$ acts on invariant test functions
having the inverse central character.

\begin{theorem}[Division-side test-function transfer]\label{III:thm:test-transfer}
Assume the specified cover of Chapter~\ref{ch:I} and Proposition~\ref{III:prop:good}.
Fix a genuine central character $\widetilde\theta$ and put
$\theta(a)=\widetilde\theta(a^*)$. The following maps are well defined:
\begin{align}
\mathcal O_{D,\widetilde\theta}:\mathcal C_{\theta^{-1}}(H)
&\longrightarrow\mathcal C_{\widetilde\theta^{-1}}(\tH),
& \mathcal O_{D,\widetilde\theta}f
&=L_D^{(\widetilde\theta^{-1})}(\mathcal A_Hf),
\label{III:eq:test-transfer-definition}\\
\mathcal B_{\widetilde\theta}:\mathcal C_{\widetilde\theta^{-1}}(\tH)
&\longrightarrow\mathcal C_{\theta^{-1}}(H),
& \mathcal B_{\widetilde\theta}\widetilde f
&=(L_D^{(\widetilde\theta^{-1})})^{-1}
  (\mathcal A_{\tH}\widetilde f).
\label{III:eq:back-transfer-definition}
\end{align}
Their images consist of conjugation-invariant test functions; and
\[
\mathcal B_{\widetilde\theta}\mathcal O_{D,\widetilde\theta}
=\mathcal A_H,\qquad
\mathcal O_{D,\widetilde\theta}\mathcal B_{\widetilde\theta}
=\mathcal A_{\tH}.
\]
They therefore induce mutually inverse linear bijections on the
respective test-function spaces modulo the kernels of their
conjugation-averaging (equivalently, normalized orbital-integral) maps.

More explicitly, for every regular $t\in T$ and $z\in\tZ$,
\begin{equation}\label{III:eq:OI-transfer}
\boxed{\quad
O_{\tH}(z\,t^*,\mathcal O_{D,\widetilde\theta}f)
=b_D\,\widetilde\theta(z)^{-1}
\frac{\Delta(t)}{\Delta(t^n)}\,O_H(t,f),
\qquad b_D=|n|_F^{-(r-1)/2}.
\quad}
\end{equation}
Consequently the normalized orbital integrals satisfy
\begin{equation}\label{III:eq:JO-transfer}
\boxed{\quad
J_{\tH}(z\,t^*,\mathcal O_{D,\widetilde\theta}f)
=b_D\,\widetilde\theta(z)^{-1}J_H(t,f).
\quad}
\end{equation}
At nongood regular lifts the genuine orbital-integral distribution
vanishes. Conversely every genuine test function $\widetilde f$ has the
same regular orbital integrals as
$\mathcal O_{D,\widetilde\theta}(\mathcal B_{\widetilde\theta}\widetilde f)$.
\end{theorem}
\begin{proof}
Conjugation averaging is a projection onto the invariant locally constant
functions, and \eqref{III:eq:ordinary-O}--\eqref{III:eq:covered-O} identify the
averages with the stated orbital integrals. The local smoothness and
invertibility of $L_D$ on invariant test functions follow from
Lemma~\ref{III:lem:well-defined} and the proof of
Theorem~\ref{III:thm:unitary}: on each regular torus quotient the star map is
a homeomorphism from $T/Z$ onto the open subgroup $(T/Z)^n$, and near
scalar elements it is a local analytic isomorphism. Thus both displayed
maps take smooth test functions to smooth test functions.
Since the image of $\mathcal A_H$ consists of invariant functions,
$\mathcal A_{\tH}\mathcal O_{D,\widetilde\theta}f
=\mathcal O_{D,\widetilde\theta}f$; applying the inverse $L_D$
gives $\mathcal B_{\widetilde\theta}\mathcal O_{D,\widetilde\theta}f
=\mathcal A_Hf$. The other composition is identical.
At a good point $z t^*$, the definition of $L_D$ applied to
$\mathcal A_Hf$ gives \eqref{III:eq:OI-transfer}. Multiplying it by
$\Delta(t^n)$ gives \eqref{III:eq:JO-transfer}. Remark~\ref{III:rem:good-vanish}
proves the vanishing on nongood lifts. The inverse composition identity
gives the converse statement. No assertion about irreducible-character
integrality is used.
\end{proof}

\begin{remark}[Comparison with the root transport of Chapter~\ref{ch:II}]\label{III:rem:root-comparison}
The orbital-integral operator $\mathcal O_{D,\widetilde\theta}$ is
\emph{not} the local root-section transport denoted $\mathcal T$
in Chapter~\ref{ch:II}. They have distinct domains, definitions and central
character conventions.  Near the center, the parameters satisfy
$t^*=\sigma(t^n)$, with $\sigma$ the root section in Chapter~\ref{ch:II};
this identity does not equate the operators on arbitrary test
functions. The orbital transfer here also involves conjugation
averaging and the factor $b_D\Delta(t)/\Delta(t^n)$.
\end{remark}

\subsection{The duality with character transfer}
The inverse central characters in Theorem~\ref{III:thm:test-transfer} are
forced by trace pairings, not a cosmetic convention.

\begin{proposition}[Trace-pairing adjunction]\label{III:prop:trace-pair}
For every smooth invariant class function $\Phi$ on $H$ of central
character $\theta$, and every $f\in\mathcal C_{\theta^{-1}}(H)$,
\begin{equation}\label{III:eq:trace-pairing}
\boxed{\quad
\int_{\tH/\tZ}
(L_D^{(\widetilde\theta)}\Phi)(\widetilde h)
 (\mathcal O_{D,\widetilde\theta}f)(\widetilde h)\,d\widetilde h
=\int_{H/Z}\Phi(h)f(h)\,dh.
\quad}
\end{equation}
In particular, for $\rho\in\Irr_\theta(H)$, writing the finite
expansion of Corollary~\ref{III:cor:matrix} as
$L_D\Theta_\rho=\sum_\tau a_{\rho\tau}\Theta_\tau$, one has
\begin{equation}\label{III:eq:character-matrix-traces}
\sum_\tau a_{\rho\tau}
\operatorname{tr}\tau(\mathcal O_{D,\widetilde\theta}f)
=\operatorname{tr}\rho(f).
\end{equation}
If $p\nmid n$, Theorem~\ref{III:thm:tame-transfer} reduces this equality to
$\operatorname{tr}\tau_\rho(\mathcal O_{D,\widetilde\theta}f)
=\operatorname{tr}\rho(f)$ for the corresponding genuine irreducible
representation $\tau_\rho$. For $p\mid n$, formula
\eqref{III:eq:character-matrix-traces} is a distributional identity with
possibly nonintegral coefficients, \emph{not} a representation
correspondence.
\end{proposition}
\begin{proof}
Because $\Phi$ is invariant, averaging $f$ does not change the integral
on the right. For the left side note that
$L_D^{(\widetilde\theta^{-1})}(\mathcal A_H f)$ is the complex
conjugate of
$L_D^{(\widetilde\theta)}(\overline{\mathcal A_H f})$:
the coefficient $b_D\Delta(t)/\Delta(t^n)$ is positive real, and
complex conjugation interchanges the two central characters.
The unitarity identity of Theorem~\ref{III:thm:unitary}, applied to $\Phi$
and $\overline{\mathcal A_H f}$, proves
\eqref{III:eq:trace-pairing}.
Equivalently, the Weyl integration formula gives the identity directly:
on each $T/Z$, the product of the two transforms contributes
$b_D^2\Delta(t)^2/\Delta(t^n)^2$ and the change of variables
\eqref{III:eq:changevar} contributes $I^{-1}$; since $b_D^2=I$, the
Weyl integrals are equal. Both trace identities follow by integrating
characters against test functions with inverse central character.
In particular \eqref{III:eq:character-matrix-traces} is a reformulation
of unitarity and surjectivity of $\mathcal O_{D,\widetilde\theta}$,
not an independent constraint forcing integrality of the matrix
$(a_{\rho\tau})$.
\end{proof}

\begin{corollary}[The identity-element factor]\label{III:cor:identity-factor}
With the probability measures and central-character conventions of this
section, the transferred invariant representative satisfies
\begin{equation}\label{III:eq:identity-factor}
\boxed{\quad(\mathcal O_{D,\widetilde\theta}f)(1)
=|n|_F^{-(r^2-1)/2}\,f(1).\quad}
\end{equation}
The factor equals $1$ at tame places. This is a local pointwise test-function
normalization; it is \emph{not} a calculation of the global identity term
in a trace formula with different Haar measures and central quotients.
\end{corollary}
\begin{proof}
The central element $1$ is fixed by conjugation, so
$(\mathcal A_Hf)(1)=f(1)$. Evaluate the continuous extension of the
transfer at $1$, using Lemma~\ref{III:lem:ratio} and
$b_D=|n|_F^{-(r-1)/2}$.
\end{proof}

\begin{remark}[The running wild-place example]
For $r=2$, $n=3$ and $F=\mathbb Q_3(\zeta_3)$,
$|3|_F=1/9$ and $b_D=3$. Thus
$J_{\tH}(z t^*,\mathcal O_{D,\widetilde\theta}f)=3\widetilde\theta(z)^{-1}J_H(t,f)$,
whereas $(\mathcal O_{D,\widetilde\theta}f)(1)=27f(1)$. This realizes the two
normalizations discussed in Remark~\ref{III:rem:wild-test}.
\end{remark}

\begin{remark}[Scope of local transfer]\label{III:rem:OI-scope}
The existence of such a transfer is elementary on a group compact
modulo its center: conjugation averaging identifies orbital integrals
with invariant function values. The content of the theorem is the
\emph{specified normalization}, extension across singular central
elements, and the trace-pairing relation. The local coefficient
$b_D$ does not determine adelic orbital measures. The distinction from Chapter~\ref{ch:II}'s root transport is made precise in Remark~\ref{III:rem:root-comparison}.
\end{remark}


\section{Kottwitz's framework for rational orbits and inner forms}
\label{III:sec:kottwitz}
This section is expository; it depends on nothing earlier in the chapter, and its terminology (stable and rational conjugacy, inner forms, the torus obstruction, Tamagawa factors) is collected in \S\ref{III:subsec:glossIII}.
Kottwitz's \emph{Stable trace formula: cuspidal tempered terms}
\cite{Kot84} explains how Galois cohomology organizes rational
conjugacy classes within stable classes (\S\S7.7--7.9), how maximal
tori are transferred to inner forms (\S9.5), and how Tamagawa
numbers enter stable trace-formula coefficients (\S\S5,8).
These are statements about \emph{linear} reductive groups; they
must not be taken as a metaplectic trace formula. In particular,
the stabilization of cuspidal tempered terms sketched in his \S12
is conditional, whereas the cohomological calculations used below
are independent of that sketch.

\paragraph{Stable and rational classes.}
For a regular semisimple element $\gamma\in G(K)$ with torus
centralizer $T=G_\gamma$, the rational conjugacy classes within
its represented stable class are parametrized by
\begin{equation}\label{III:eq:kottwitz-rational}
\ker\!\left[H^1(K,T)\longrightarrow H^1(K,G)\right]
\end{equation}
\cite[\S7.7]{Kot84}. In the present setting,
$G=\GL_r$ or $G'=\GL_1(D)$ and
$T=\operatorname{Res}_{E/K}\mathbb G_m$ for a maximal subfield
$E/K$ of degree $r$. By Shapiro's lemma and Hilbert~90,
$H^1(K,T)=H^1(E,\mathbb G_m)=1$.
Thus every \emph{represented} regular stable class comprises
exactly one rational conjugacy class. This does not decide
whether the class has a good lift in a covering group.

\paragraph{The inner-form embedding criterion.}
A degree-$r$ extension $E/K$ embeds in $D$ as a maximal
subfield if and only if
\begin{equation}\label{III:eq:brauer-embedding}
[D\otimes_K E]=0\quad\text{in }\operatorname{Br}(E).
\end{equation}
By Brauer--Hasse--Noether, this condition can be checked at all
completions. Where $D_v$ is division of degree $r$, the local
algebra $E\otimes_K K_v$ must be a field of degree $r$.
Consequently, for these inner forms, local embeddability at every
place implies global embeddability. Kottwitz's general \S9.5
formulates an obstruction to torus transfer in a class of
$H^2(K,T_{\mathrm{sc}})$ (pp.~638--639); the specialized Brauer
criterion avoids that additional complication here. For linear
orbital-integral transfer see also \cite{Coh18}.

\paragraph{Tamagawa factors.}
Kottwitz \S5 identifies the relative Tamagawa number
$\tau_1(G)=\tau(G)/\tau(G_{\mathrm{sc}})$ as
\begin{equation}\label{III:eq:kottwitz-tamagawa}
\tau_1(G)=
\frac{|\pi_0(Z(\widehat G)^{\Gamma_K})|}
{|\ker^1(K,Z(\widehat G))|}.
\end{equation}
For $\GL_r$ and $\GL_1(D)$ the dual center is
$\mathbb C^\times$ with trivial Galois action, so this expression
is $1$. Kottwitz \S8 uses relative Tamagawa numbers in stable
trace-formula coefficients. Our local $b_D$ instead uses
\emph{probability} Haar measures on compact quotients; its
conversion to a coherent adelic convention remains a separate
calculation, not a consequence of \eqref{III:eq:kottwitz-tamagawa}.

\section{Global power classes and a Hasse principle}\label{III:sec:global-power}
The scalar-modulo-$n$th-power quotient is typically infinite,
but its subgroup of classes trivial outside a finite set of places
vanishes under the hypotheses below. The proof separates the
norm argument from the Kummer--Chebotarev detection step. The section depends on \S\ref{III:sec:power} and on global class field theory; it introduces the scalar power quotient $\mathcal Q(E/K,n)$ and the group $\Sha^1(K,T[n])$.
The following large quotient measures the failure of \emph{global
surjectivity} of a power map, but its size is not the correct measure
of failure of a \emph{local-to-global principle}. Let $E/K$ be a maximal subfield
of $D$, of degree $r$, and consider the multiplicative quotient
\begin{equation}\label{III:eq:rational-obstruction}
\mathcal Q(E/K,n)=E^\times/(K^\times E^{\times n}).
\end{equation}
A regular rational class represented by $\gamma\in E^\times$
can lie in the scalar multiple of the rational $n$th-power image
only if its class in \eqref{III:eq:rational-obstruction} is trivial.
This obstruction must not be confused with the finite local index
of Lemma~\ref{III:lem:index}.

\begin{proposition}[The ambient power quotient is infinite]
\label{III:prop:global-cokernel}
Let $E/K$ be a number-field extension of degree at least two and
let $n>1$. Then $\mathcal Q(E/K,n)$ is infinite. More precisely,
for any $m\geq 1$ and any distinct finite places
$\mathfrak p_1,\dots,\mathfrak p_m$ of $K$ splitting completely in
$E$, there is a surjective homomorphism
\[
\mathcal Q(E/K,n)\longrightarrow (\mathbb Z/n\mathbb Z)^m.
\]
In particular, even when $\mu_n\subset K$ and $\gcd(n,r)=1$,
the map $[t]\mapsto[t^n]$ on $E^\times/K^\times$ cannot be
surjective.
\end{proposition}
\begin{proof}
For each $i$ choose two distinct primes
$\mathfrak P_{i,1}$ and $\mathfrak P_{i,2}$ of $E$ above
$\mathfrak p_i$. Because $\mathfrak p_i$ splits completely,
for $a\in K^\times$ the normalized discrete valuations satisfy
$v_{\mathfrak P_{i,1}}(a)=v_{\mathfrak P_{i,2}}(a)$.
Consequently
\[
\partial_i([x])=
 v_{\mathfrak P_{i,1}}(x)-v_{\mathfrak P_{i,2}}(x)
 \pmod n
\]
is well defined on $\mathcal Q(E/K,n)$: multiplication by
$K^\times$ changes neither difference, and multiplication by
an $n$th power changes it by a multiple of $n$.
The weak approximation theorem for $E$ permits independent
prescription of the valuations at this finite collection of primes,
so the product $(\partial_1,\dots,\partial_m)$ is surjective.
Chebotarev's theorem supplies infinitely many completely split
primes; thus $m$ is arbitrary and the quotient is infinite.
\end{proof}

\begin{remark}[Effect on the proposed trace formula]
The infinitude of $\mathcal Q(E/K,n)$ does \emph{not} show that
there is an infinite global obstruction among the classes satisfying
all local power conditions. Indeed, the classes exhibited by valuation
differences are already nontrivial in the local power quotient at
the corresponding completely split places. For the split-cover
commutator these may also be nongood; that additional implication
must be checked in the particular Flicker--Kazhdan cocycle, not
inferred merely from the valuation calculation. One must distinguish
a large ambient quotient from the kernel of the localization map.
That kernel is controlled by the next proposition.
\end{remark}


\begin{proposition}[Local-to-global powers modulo scalars]\label{III:prop:global-power-hasse}
Let $E/K$ be a separable extension of number fields of degree $r$,
let $n>1$, and assume $\gcd(n,r)=1$ and $\mu_n\subset K$.
For \emph{every finite set of places} $\Sigma$ of $K$, the natural map
\begin{equation}\label{III:eq:power-local-inj}
 \mathcal Q(E/K,n)\longrightarrow
 \prod_{v\notin\Sigma}\frac{(E\otimes_K K_v)^\times}
 {K_v^\times (E\otimes_K K_v)^{\times n}}
\end{equation}
is injective. Thus being a scalar times an $n$th power at
\emph{almost all places} already implies being so globally.
\end{proposition}
\begin{proof}
Let $t\in E^\times$ map to the trivial class outside a finite set $\Sigma$.
Choose integers $u,k$ satisfying $ru+nk=1$, and set
$a=N_{E/K}(t)^u\in K^\times$.
For each $v\notin\Sigma$, the assumption supplies
$t=a_v y_v^n$ in $(E\otimes_K K_v)^\times$, for
$a_v\in K_v^\times$ and $y_v\in(E\otimes_K K_v)^\times$.
Taking the norm from this degree-$r$ \emph{\'{e}tale algebra}
produces $N(t)=a_v^r N(y_v)^n$. Consequently
\[
\frac{t}{a}=
\left(a_v^k\frac{y_v}{N(y_v)^u}\right)^n
\quad\text{in }(E\otimes_K K_v)^\times.
\]
Thus $t/a$ is an $n$th power at each completion $E_w$ lying over a place $v\notin\Sigma$.
Since $\mu_n\subset E$, a nontrivial Kummer extension
$E((t/a)^{1/n})/E$ would be a nontrivial cyclic Galois
extension and, by Chebotarev, would possess infinitely many unramified
places with nontrivial Frobenius. We may choose one lying over
a place $v\notin\Sigma$. At that completion $t/a$ could not
be an $n$th power, a contradiction. Hence
$t/a\in E^{\times n}$ and $[t]=1$ in $\mathcal Q(E/K,n)$.
\end{proof}

\begin{corollary}[Vanishing of the associated \texorpdfstring{$\Sha^1$}{Sha-1}]\label{III:cor:sha-power}
Under the assumptions of Proposition~\ref{III:prop:global-power-hasse},
put $T=\operatorname{Res}_{E/K}\mathbb G_m/\mathbb G_m$ and
let $T[n]$ be its finite $n$-torsion group scheme. Then for every finite $\Sigma$, the restricted localization kernel
\[\Sha^1_{\Sigma}(K,T[n]):=\ker\bigl(H^1(K,T[n])\to\prod_{v\notin\Sigma}H^1(K_v,T[n])\bigr)\]
vanishes. In particular, the usual
\begin{equation}\label{III:eq:sha-zero}
\Sha^1(K,T[n]):=
\ker\left[H^1(K,T[n])\longrightarrow
                \prod_v H^1(K_v,T[n])\right]=0.
\end{equation}
This is stronger than the general
finiteness of $\Sha^1$ for a finite Galois module
\cite[Chapter~I, \S4]{MilADT06}.
\end{corollary}
\begin{proof}
From $1\to\mathbb G_m\to\operatorname{Res}_{E/K}\mathbb G_m
\to T\to1$, Hilbert~90 gives $T(K)=E^\times/K^\times$.
The same cohomology sequence gives
\[
 H^1(K,T)\cong
 \ker\bigl(\operatorname{Br}(K)\to\operatorname{Br}(E)\bigr).
\]
Restriction followed by corestriction is multiplication by $r$,
so this group is killed by $r$. Hence, as $\gcd(n,r)=1$,
$H^1(K,T)[n]=0$. The Kummer exact sequence for multiplication
by $n$ on $T$ therefore gives
\[
 H^1(K,T[n])\cong T(K)/nT(K)=\mathcal Q(E/K,n).
\]
At each $K_v$ the same argument applies to the finite \'{e}tale
algebra $E\otimes_K K_v$ (its total degree is still $r$), yielding
analogous isomorphisms for $H^1(K_v,T[n])$.
Under these identifications the localization map is
\eqref{III:eq:power-local-inj}, which is injective, even after omitting any finite set of places, by the proposition.
\end{proof}

\begin{remark}[Scope of the Hasse principle]\label{III:rem:sha-limits}
The conclusion uses local power conditions at almost all places,
not merely at the division places $S$. Moreover, an $n$th-power
condition in the scalar quotient must not be identified with the
good-element condition for the cover of a split group without a
separate commutator calculation. The theorem eliminates this
particular Kummer obstruction; it does not provide the missing
metaplectic geometric comparison.
\end{remark}



\section{The global comparison: remaining problems}\label{III:sec:global}
The following spectral identity is a target rather than a theorem.
The arithmetic in Section~\ref{III:sec:global-power} simplifies its
rational-orbit accounting, but leaves the covering-group good-locus
and measure comparisons untouched. The section depends on \S\S\ref{III:sec:orbital}--\ref{III:sec:global-power} and introduces the target identity \eqref{III:eq:global-conj} and the map $\Psi_D$.
Let $K$ be a number field containing $\mu_n$ (for odd $n>1$ this forces $K$ to be totally imaginary, as in Chapters~\ref{ch:I} and~\ref{ch:II}), and $D/K$ a central division algebra of prime degree $r$. Denote by $S$ its finite ramification set. We consider the automorphic spectra of $D^\times(\A_K)$ and its chosen metaplectic cover, both modulo their respective centers. The quotient of $D^\times(\A_K)$ by its rational points and adelic center is compact, hence both spectral trace formulas are discrete. \emph{This fact alone does not establish equality of their geometric sides.}

The following proposed identity is the main mathematical target:
\begin{equation}\label{III:eq:global-conj}
\sum_{\widetilde\Pi_D}m(\widetilde\Pi_D)
\operatorname{tr}\widetilde\Pi_D(\widetilde f)
\stackrel{?}{=}
\sum_{\Pi_D}m(\Pi_D)\operatorname{tr}\Pi_D(f),
\end{equation}
for appropriately matched factorizable test functions and compatible global central characters. It is \emph{conditional} on a comparison of the geometric sides with all constants correct. Section~\ref{III:sec:orbital} records a normalized division-side local
matching operator on test functions, dual to the character-space
transform. Its existence is automatic from compactness modulo the
center, while smoothness and the normalization require the stated
local input. Kottwitz's framework in Section~\ref{III:sec:kottwitz}
clarifies the ordinary rational-orbit and measure issues, but neither
automatically gives a geometric comparison for the cover nor
settles wild-place character integrality.

\subsection{Four proof obligations}
\begin{enumerate}[label=(\arabic*)]
\item \textbf{Rational orbit accounting.} Combine the uniqueness of represented regular rational classes from
Section~\ref{III:sec:kottwitz} with the local-to-global power result
of Proposition~\ref{III:prop:global-power-hasse}. Then determine which
rational elements are locally \emph{good} for the particular
cover at split and division places. Vanishing of the Kummer
localization kernel alone does not identify the surviving cover orbits.
\item \textbf{Local normalization and compatibility.} At $v\in S$, convert the probability-measure test transfer of Section~\ref{III:sec:orbital} to the fixed global measures and check the local geometric factor against the exact desired global comparison. At split places, cite the exact usable form of Flicker--Kazhdan \cite{FK86}, taking into account Zou's correction to their Lemma 27.1 \cite{Z23}. Check unramified unit elements and archimedean factors separately.
\item \textbf{Identity and volume terms.} With the trace-formula quotient measures fixed, calculate the identity terms, orbital-integral Jacobians, and central character factors. The $b_D$ of Theorem~\ref{III:thm:unitary} is a \emph{local probability-measure normalization} and does not dispose of the global normalization problem.
\item \textbf{Spectral separation and integrality.} Derive a finite, or otherwise rigorously summable, character identity. Use classical strong multiplicity one for the linear group \cite{Bad08}, but do not infer integrality merely from $L^2$ unitarity. The proposed Cauchy--Schwarz argument requires an explicit identity, its correct normalization, and control of every other local factor.
\end{enumerate}

\begin{remark}[Restriction on $n$]\label{III:rem:N}
The standard hypotheses under which one invokes the original Flicker--Kazhdan trace formula include $(n,N_r)=1$, where $N_r$ is the least common multiple of the composite integers $\le r$. Our local Theorem~\ref{III:thm:unitary} uses only $n$ odd, $(n,r)=1$, and the covering result in Proposition~\ref{III:prop:good}; it does not use $(n,N_r)=1$. The global comparison will provisionally impose the latter condition until the exact later improvements or alternatives are verified.
\end{remark}

\subsection{Conditional representation-theoretic consequences}
If \eqref{III:eq:global-conj} is established in a form that proves character integrality at the remaining \emph{wild} ramified places, while Theorem~\ref{III:thm:tame-transfer} handles the tame ones, the resulting map $\Psi_D$ from irreducible genuine representations of the cover to $\Irr(D_v^\times)$ can be composed, subject to central-character and normalization checks, with
\[
\Irr(D_v^\times)\xrightarrow{\ \JL\ }
\Irr_{\mathrm{ds}}(\GL_r(F))
\xleftarrow{\ \mathrm{MC}\ }
\Irr_{\mathrm{gen,ds}}(\widetilde{\GL}_r(F)),
\]
where $\mathrm{MC}$ is Zou's corrected form of the Flicker--Kazhdan correspondence, with the required restriction that the linear central character be trivial on $\mu_n$. At tame places this composition gives the proved conclusion of
Corollary~\ref{III:cor:D2tame} in the stated range of the external
Flicker--Kazhdan theorem; outside that range it remains conditional
on an applicable extension of the external input.
This would establish a local metaplectic Jacquet--Langlands correspondence. The fact that local supercuspidals are easier to isolate spectrally does not eliminate the wild-place integrality obstruction for all representations.

Similarly, removing the non-cancellation hypothesis $(H_D')$ of \cref{II:def:HDprime}, would require showing that a prescribed packet has only finitely many $S$-types, or some other condition contradicting the cancellation criterion from Chapter~\ref{ch:II}. The original planning memo's argument through strong multiplicity one is a plausible \emph{conditional} route, but its spectral separation and finiteness arguments have not yet been established.

Removal of the extra cuspidal place $w_0$, or of the converse hypothesis $(H')$ (\cref{II:cor:Hprimesq}), remains a different problem. None of these goals is asserted as a theorem in the present draft.
\endgroup
